\section{Conclusion}

The table collects the approximation factors proved for the two depth
objectives and the two answer models. Here $\gamma$ is the calibrated
count-potential coefficient from Section~\ref{sec:worst-calibration},
and $\rho_{\rm two}$ is the calibrated expected-depth factor from
Section~\ref{sec:two-phase}.

\begin{center}
\begin{tabular}{p{0.37\textwidth}p{0.24\textwidth}p{0.27\textwidth}}
\hline
Objective & Polynomial time & Quasipolynomial time\\
\hline
Maximum depth, binary answers &
$2+\gamma<4.464966$ & $2+\varepsilon$\\[3pt]
Maximum depth, on-line answers &
$4+2\gamma<8.929932$ & $4+\varepsilon$\\[3pt]
Expected depth, binary answers and arbitrary priors &
$\rho_{\rm two}<5.182976$ & $2+\varepsilon$\\[3pt]
Expected depth, on-line answers and arbitrary priors &
$2\rho_{\rm two}<10.365952$ & $4+\varepsilon$\\
\hline
\end{tabular}
\end{center}

The common-witness normalization is the structural step shared by
both algorithms. In a flat state, it turns each supplied test into
an opposite cap and represents that cap by at most two adjacent
vertices of a convex hull. For each hypothesis, the positive hull
answers form a circular interval. This representation keeps the
query family fixed and gives the supplied tests an interval view.

The two objectives use this shared structure in different ways. To
obtain the maximum-depth guarantee, we follow a common witness through
an optimal tree. Restricting that tree to a pattern class removes the
queries on the witness prefix, so that prefix gives the class a
deadline. The adjacent-support lemma reduces these deadlines to
singleton and pair blocks, which the scheduler orders. Endpoint
decoding then uses cached negative answers to preserve the remaining
height credit. When there are logarithmically many discovery slots, an
exact Cartesian-tree dynamic program improves the schedule; the global
cap and priority fallback handle the other cases. Together these steps
give a factor below $4.464966$.

For expected depth, the algorithm first learns each full hull pattern
and then recurses within its class. Restricting an optimal tree to a
class removes queries whose answers are already determined, and that
saved cost helps pay for the first positive hull answer. The answers
also reveal information about the pattern, which shortens the endpoint
code. Exact integer-budget coverage over adjacent hull blocks selects
when each pattern is found, while the all-negative waiting charge
accounts for the zero-pattern route. The schedule search chooses the
first two budgets, geometric shift, and endpoint mode together, because
asking for a second endpoint can reduce later decoding work. The signed
first-cover correction registers how moving a batch boundary changes
each class's first-hit rank, while the endpoint term credits the
information a second answer provides. Combined with restriction
savings, these credits telescope over the recursive tree, so each stage
cost is charged once.

The quasipolynomial algorithms store compatible discovery grades.
This lets them handle all deadlines together. Their internal bounds
have leading coefficient arbitrarily close to two and an additive
count or entropy potential; exact logarithmic-depth lookahead absorbs
that potential and gives the final $(2+\varepsilon)$ ratio.
The on-line-answer entries follow from the same factor-two simulation,
using the binary guarantee with parameter $\varepsilon/2$.
The algorithms run in quasipolynomial time.

For randomized per-hypothesis depth, the adversary chooses the hidden
hypothesis before sampling the random tree. Multiplicative weights
then constructs trees for a sequence of positive rational priors.
A restricted mixture LP combines these trees to minimize the largest
per-hypothesis expected depth. Proposition~\ref{prop:randomized-duality}
identifies the hardest prior as the comparison benchmark, and
Theorem~\ref{thm:randomized-transfer} shows that the weight updates
preserve the expected-depth oracle's approximation factor up to an
arbitrary tolerance. The LP compresses the resulting strategy to at
most $n$ trees. Since this objective differs from deterministic
maximum depth, the LP value and its comparison priors provide a
separate, instance-specific approximation certificate.

The small-cap results also delimit possible hardness constructions.
With a distinguished all-miss hypothesis and fixed cap size, maximum
depth is exactly solvable in polynomial time. For expected depth, the
weighted singleton/pair path-and-cycle subclass with all singleton
tests has an EPTAS. Consequently, APX-hardness for the current weighted
pair-path reduction would imply $\mathrm P=\mathrm{NP}$. A broader
uniform subclass also has a PTAS: fixed-size caps in convex position
with an all-miss hypothesis, even when singleton tests are absent.
Clearing short gaps permits exact optimization of the remaining block
interleavings, and the prefix cost stays a small fraction of the full
optimum. A constant-factor hardness construction must therefore use
instances outside this uniform fixed-cap class.

The hardness reductions show that strong NP-completeness persists even
when all hypotheses lie on a parabola and every query is a supplied
line. The expected-depth reduction applies to arbitrary positive priors
and to uniform priors: replacing each weight by a uniform cloud adds a
fixed internal identification cost to every matching. Several
approximation questions remain: whether constant-factor
inapproximability holds, whether the polynomial factors can be lowered,
and whether the quasipolynomial factors can fall below two.
